\documentclass[conference]{IEEEtran}
\IEEEoverridecommandlockouts
\usepackage{amsmath,amssymb}

\setlength{\parskip}{0pt}

\begin{document}

\title{Riptide: Breaking the Sequential Barrier in DEFLATE Decompression}

\author{
\begin{tabular}{@{}c@{\hspace{2.2cm}}c@{}}
\textbf{Devansh Hiren Shah} & \textbf{Mantri Samarth Prakash}\\
Dept. of Information Technology & Dept. of Information Technology\\
NIT Karnataka, Surathkal --- 241IT024 & NIT Karnataka, Surathkal --- 241IT043\\[8pt]
\textbf{Shubh Minesh Thakkar} & \textbf{Ritesh Bhagat}\\
Dept. of Information Technology & Dept. of Information Technology\\
NIT Karnataka, Surathkal --- 241IT076 & NIT Karnataka, Surathkal --- 241IT063
\end{tabular}
}

\maketitle

\section{Introduction}

DEFLATE is the compression format behind \texttt{gzip}, ZIP, PNG, HTTP
transfers and the default codecs of Apache Parquet and ORC. Compressing with it
parallelises easily, since the compressor may cut the input wherever it likes.
Decompressing a single \texttt{.gz} file does not: no practical decoder turns
extra CPU cores into extra speed on one stream. Big-data engines can split a
file across tasks only if a worker can begin reading from the middle. Gzip
forbids this, so a 10\,GB \texttt{.gz} file becomes one task on one core even
on a 500-core Spark cluster. The same ceiling appears elsewhere: genomics
pipelines align reads at tens of gigabytes per second yet are throttled by one
thread decompressing \texttt{.fastq.gz}. As data consumers grow more parallel,
this gap widens with every new machine. The standard advice, to re-encode data
into a splittable format, first requires decompressing everything sequentially,
which is the very cost we wish to remove.

Given an unmodified gzip stream and $P$ processors, we therefore aim to produce
output bit-identical to the standard sequential decoder, in time that shrinks
as $P$ grows, without re-encoding the file and without an index built by a
prior sequential pass. Three obstacles stand in the way. First, the stream
holds no directory of block positions, and each block's start is found only
after parsing the previous one, so the split points we need are exactly what
the sequential scan exists to compute. Second, Huffman symbols vary in length
and are packed without byte alignment, so a worker starting at a wrong bit
position does not fail loudly but silently emits wrong output. Third, a match
instructs the decoder to copy bytes from up to 32\,KB earlier in the output, so
one chunk can genuinely need data another has not yet produced.

Existing tools avoid these obstacles rather than solve them. \texttt{pigz}
overlaps only file I/O when decompressing, leaving the decode loop
single-threaded, while Intel's \texttt{igzip} uses SIMD to make one thread
faster without scaling with core count. Chunked formats such as BGZF and
index-based tools such as \texttt{gztool} work only if the file is rewritten in
a different layout or scanned sequentially first, so neither helps a first-time
read of an ordinary \texttt{.gz}. \texttt{rapidgzip} is the strongest prior
result, genuinely decoding standard gzip in parallel, but it is CPU-only and
its recovery pass inflates memory, and NVIDIA's nvCOMP reaches GPU speeds only
inside its own chunked container. The gap is that no decoder handles
unmodified, unindexed gzip with one strategy running on both multicore CPUs and
GPUs.

Such a decoder would accelerate data that already exists, without migration or
re-encoding: archived genomic reads, Spark and Trino scans over
gzip-compressed Parquet, queries over gzipped logs, and GPU data loading, where
decoding on the device also avoids sending the larger uncompressed form across
PCIe. The beneficiaries are bioinformatics researchers, data-engineering teams,
and machine-learning infrastructure operators.

Our approach rests on the observation that the dependence is weaker than it
appears. Reading compressed bits into symbols needs no previously decoded byte;
only executing a copy does. Riptide thus parses in parallel and records any
copy reaching before its own chunk as a placeholder, resolving all placeholders
in a second parallel pass. Because a match looks back at most 32\,KB, only the
first 32\,KB of each chunk can depend on another, so with large chunks well
over 99\% of the work is dependence-free. To our knowledge this is the first
design combining index-free discovery of block boundaries, a validation step
that proves rather than assumes a genuine symbol boundary, and
placeholder-based deferred copying, in one formulation intended for both CPUs
and GPUs. The checksum already stored in every gzip file then verifies the
entire speculative result at no extra cost. Our underlying claim is that
DEFLATE is sequential because of how decoders are written, not because the
format demands it.

We will evaluate correctness by requiring byte-for-byte agreement with
\texttt{zlib} across all block types, unusual matches and corrupted files.
Performance will be measured as speedup against thread count and as throughput
compared with \texttt{zlib}, \texttt{igzip}, \texttt{rapidgzip} and nvCOMP,
alongside peak memory and sensitivity to compression ratio.

We expect to contribute algorithms for boundary discovery, verified entry
points and deferred copy resolution; Riptide, a decompressor with CPU and GPU
backends for unmodified \texttt{.gz} files; and measurements of how much
parallelism real gzip files contain.

\end{document}
